\documentclass[12pt,a4paper]{article}
\usepackage[utf8]{inputenc}
\usepackage[T1]{fontenc}
\usepackage{amsmath,amssymb,amsthm}
\usepackage[margin=1in]{geometry}
\usepackage{hyperref}
\usepackage{natbib}
\usepackage{enumitem}

% Theorem environments
\newtheorem{definition}{Definition}[section]
\newtheorem{lemma}{Lemma}[section]
\newtheorem{theorem}{Theorem}[section]
\newtheorem{corollary}{Corollary}[section]
\newtheorem{proposition}{Proposition}[section]
\newtheorem*{remark}{Remark}

% RA-specific commands
\newcommand{\Pact}{P_{\mathrm{act}}}
\newcommand{\Tact}{T^{\mu\nu}_{\mathrm{act}}}
\newcommand{\Srel}{\Delta S(\rho\|\sigma_0)}
\newcommand{\DAG}{\mathcal{G}}

\title{Actualized Geometry and a Null Prediction\\
for Gravity-Mediated Entanglement\\[6pt]
{\large The Bose--Marletto--Vedral Protocol in Relational Actualism}}

\author{Joshua F.\ Sandeman\\
{\small Independent Researcher, Salem, Oregon}\\
{\small \texttt{sansuikyo@gmail.com}}\\
{\small \texttt{relationalactualism.org}}\\[6pt]
{\small with Claude (Anthropic) and ChatGPT (OpenAI)}}

\date{April 2026}

\begin{document}
\maketitle

\begin{abstract}
We show that Relational Actualism (RA) predicts a strict null result
for gravity-mediated entanglement in the Bose--Marletto--Vedral (BMV)
protocol. The prediction follows from a single architectural principle:
the effective spacetime metric is sourced only by the actualized causal
ledger---the growing directed acyclic graph of irreversible physical
events---not by unactualized quantum branch structure. During coherent
spatial superposition, the positional alternatives of the test masses
remain in the reversible possibility layer (Step~2 of the Engine of
Becoming) and have not yet been inscribed as actualized vertices
(Step~4). Since the metric is recalculated only from actualized
vertices (Step~5), no branch-resolved geometry exists during the
coherent evolution interval. The gravitational phase is therefore
common-mode across all positional branches, and the entangling
invariant vanishes identically: $\Delta\phi_{\mathrm{grav}} = 0$.
We prove a coherence--geometry incompatibility theorem showing that
the two conditions required by the standard BMV mechanism---maintained
positional coherence and branch-resolved gravitational geometry---are
structurally mutually exclusive in RA. A positive BMV signal, after
exclusion of nongravitational channels, would directly falsify the
Step-2/Step-5 separation that is central to the RA architecture.
\end{abstract}

\tableofcontents

%% ===============================================================
\section{Introduction}
\label{sec:intro}
%% ===============================================================

The Bose--Marletto--Vedral (BMV) protocol \cite{Bose2017,Marletto2017}
proposes to test whether gravity can mediate quantum entanglement
between two mesoscopic masses, each placed in a coherent spatial
superposition. If gravity is fundamentally quantum, the branch-dependent
Newtonian interaction generates a nonseparable phase that entangles
the masses. If gravity is fundamentally classical or---as Relational
Actualism (RA) holds---fundamentally \emph{actualized}, no such
entanglement arises.

The standard BMV argument implicitly assumes that each branch of the
spatial superposition sources its own gravitational field prior to
any measurement or actualization event. In the language of RA, this
amounts to assuming that the effective metric responds to the Level-2
quantum-measure amplitude structure---the space of unrealized
possibilities---rather than exclusively to the Level-3 actualized
causal ledger.

RA denies precisely this assumption. The framework's central
architectural feature is a sharp separation between three levels of
physical description:
\begin{itemize}[nosep]
\item \textbf{Step 2 (Propose):} Reversible quantum amplitudes are
assigned to candidate continuations of the actualized graph. No
vertices are written.
\item \textbf{Step 4 (Inscribe):} One candidate is selected and
irreversibly written into the causal DAG as a new actualized vertex,
subject to the BDG filter $S_{\mathrm{BDG}} > 0$.
\item \textbf{Step 5 (Update geometry):} The effective metric
$g_{\mu\nu}$ is recalculated from the updated actualized graph via
the coarse-graining map $g_{\mu\nu} = \mathcal{G}[G_n]$.
\end{itemize}

\noindent
The BMV protocol requires Step~5 to respond at Step~2. RA says Step~5
responds only at Step~4. That is the entire content of the null
prediction, and this note makes it precise.

\medskip
\noindent
\textbf{This paper does and does not:}
\begin{itemize}[nosep]
\item[\checkmark] Proves a null BMV theorem from RA's Step-2/Step-5
separation.
\item[\checkmark] States the precise falsification condition for RA.
\item[\checkmark] Distinguishes RA's prediction from both quantized
gravity and semiclassical gravity.
\item[$\times$] Does not derive the full RA ontology (see companion
papers \cite{RAQM,RAGC,EB}).
\item[$\times$] Does not claim gravity is ``classical'' in the naive
sense.
\end{itemize}

%% ===============================================================
\section{The Standard BMV Mechanism}
\label{sec:standard}
%% ===============================================================

Two equal masses $A$ and $B$, each of mass $m$, are prepared in the
spatial superposition
\begin{equation}
|{+}\rangle_X = \frac{1}{\sqrt{2}}
\bigl(|L\rangle_X + |R\rangle_X\bigr),
\qquad X \in \{A,B\}.
\end{equation}
The initial joint state is
\begin{equation}
|\Psi_0\rangle = |{+}\rangle_A \otimes |{+}\rangle_B
= \frac{1}{2}\bigl(|LL\rangle + |LR\rangle + |RL\rangle
+ |RR\rangle\bigr).
\end{equation}
Each branch $|ij\rangle$ has separation $d_{ij}$ and Newtonian
interaction energy $U_{ij} = -Gm^2/d_{ij}$. After interaction
time~$t$, branch $|ij\rangle$ accumulates the gravitational phase
\begin{equation}
\phi_{ij} = \frac{Gm^2 t}{\hbar\, d_{ij}}.
\end{equation}
The evolved state is
\begin{equation}
|\Psi_{\mathrm{std}}(t)\rangle
= \frac{1}{2}\sum_{i,j \in \{L,R\}} e^{i\phi_{ij}} |ij\rangle.
\end{equation}
This state is entangled whenever the entangling phase invariant
\begin{equation}
\Delta\phi := \phi_{LL} + \phi_{RR} - \phi_{LR} - \phi_{RL}
\end{equation}
is nonzero modulo $2\pi$. For generic separations,
$\Delta\phi \neq 0$, and the standard prediction is positive
gravity-mediated entanglement.

The implicit assumption behind this calculation is that each branch
$|ij\rangle$ has its own gravitational potential, sourced by the
mass configuration that branch represents. That is the assumption
RA removes.

%% ===============================================================
\section{RA Architecture: The Step-2/Step-5 Separation}
\label{sec:architecture}
%% ===============================================================

In Relational Actualism, the universe is represented by a growing
causal directed acyclic graph (DAG)~$G = (V, \prec)$ whose vertices
are actualized events and whose edges encode causal precedence. The
dynamics proceeds through the five-step Engine of Becoming:

\begin{enumerate}[nosep]
\item \textbf{Survey:} Identify the boundary of $G$ (vertices with
no future-directed edges).
\item \textbf{Propose:} Assign complex amplitudes to candidate
extensions of $G$, weighted by the BDG action
$S = 1 - N_1 + 9N_2 - 16N_3 + 8N_4$.
\item \textbf{Filter:} Retain only candidates with
$S_{\mathrm{BDG}} > 0$.
\item \textbf{Inscribe:} Select one candidate and write it
irreversibly into the graph as a new vertex.
\item \textbf{Update:} Recalculate the effective metric from the
updated actualized graph: $g_{\mu\nu} = \mathcal{G}[G]$.
\end{enumerate}

The critical architectural feature for BMV is the separation between
Steps~2 and~5. Step~2 creates a weighted space of possibilities.
Step~5 responds only to what has been inscribed at Step~4.

\begin{definition}[Ledger-sourced metric]
\label{def:ledger}
Let $G_n = (V_n, E_n)$ be the actualized causal DAG after $n$
actualization events. The emergent metric at stage~$n$ is
\begin{equation}
g_{\mu\nu}^{(n)} = \mathcal{G}[G_n],
\end{equation}
where $\mathcal{G}$ is the Step-5 coarse-graining map. This metric
depends only on actualized vertices and edges, not on unactualized
candidate extensions present in the Step-2 amplitude layer.
\end{definition}

\begin{definition}[Branch-coherent regime]
\label{def:coherent}
A BMV protocol is in the \emph{branch-coherent regime} on the
interval $[0,T]$ if, for each mass, the positional alternatives
$L$ and $R$ remain in coherent superposition throughout $[0,T]$:
no irreversible positional actualization event occurs during that
interval. Equivalently, the positional degree of freedom remains at
Step~2 rather than passing through Step~4.
\end{definition}

%% ===============================================================
\section{The RA-BMV Null Theorem}
\label{sec:theorem}
%% ===============================================================

\begin{lemma}[Step-2 non-sourcing]
\label{lem:nonsource}
Let $\mathcal{C}_n$ denote the set of candidate extensions available
at Step~2 from ledger state $G_n$. Suppose a degree of freedom
remains fully within the Step-2 amplitude layer during $[0,T]$, so
that none of its alternatives are selected at Step~4 during that
interval. Then the Step-5 metric is insensitive to the decomposition
of amplitudes across those alternatives.

In symbols: if no candidate in the relevant sector actualizes, then
for that sector
\begin{equation}
G(t) = G_n \quad \text{for all } t \in [0,T],
\end{equation}
and therefore
\begin{equation}
g_{\mu\nu}(t) = \mathcal{G}[G_n] \quad \text{for all } t \in [0,T].
\end{equation}
\end{lemma}

\begin{proof}
Step~2 enumerates weighted possibilities; it does not add vertices to
the actualized graph. By hypothesis, no positional branch is selected
at Step~4 during $[0,T]$. Therefore the actualized graph in the
relevant sector does not change. Since Step~5 depends only on the
actualized graph (Definition~\ref{def:ledger}), the metric remains
unchanged throughout the interval.
\end{proof}

\begin{lemma}[No branch-indexed metric before actualization]
\label{lem:nobranch}
In the branch-coherent regime, there is no ontically meaningful
family
\begin{equation}
\{g^{LL}_{\mu\nu},\; g^{LR}_{\mu\nu},\;
g^{RL}_{\mu\nu},\; g^{RR}_{\mu\nu}\}
\end{equation}
of simultaneously available branch-resolved metrics.
\end{lemma}

\begin{proof}
Such a family would require the four branch labels $LL, LR, RL, RR$
to correspond to four distinct actual source configurations. But in
the branch-coherent regime, those labels denote only unactualized
Step-2 alternatives. By Lemma~\ref{lem:nonsource}, Step~5 has access
only to the common actual ledger, not to branch labels that have not
passed through Step~4.
\end{proof}

\begin{theorem}[RA-BMV null theorem]
\label{thm:null}
Consider a BMV protocol with two masses $A$ and $B$, each prepared
in a coherent two-position superposition. Assume:
\begin{enumerate}[nosep]
\item the protocol remains in the branch-coherent regime during the
interaction interval $[0,T]$;
\item the emergent metric is Step-5 sourced from the actualized
ledger only;
\item nongravitational entangling channels are absent.
\end{enumerate}
Then the purely gravitational contribution to the four branch phases
is common-mode:
\begin{equation}
\phi^{\mathrm{grav}}_{LL}
= \phi^{\mathrm{grav}}_{LR}
= \phi^{\mathrm{grav}}_{RL}
= \phi^{\mathrm{grav}}_{RR},
\end{equation}
and the entangling phase invariant vanishes:
\begin{equation}
\boxed{\Delta\phi_{\mathrm{grav}} := \phi_{LL} + \phi_{RR}
- \phi_{LR} - \phi_{RL} = 0.}
\end{equation}
Hence the protocol cannot generate gravity-mediated entanglement.
\end{theorem}

\begin{proof}
By Lemma~\ref{lem:nobranch}, no branch-resolved metric exists during
the coherent interval. Gravity contributes the same Step-5 background
potential $U_0 = U[g_{\mu\nu}(G_n)]$ to each amplitude component.
The induced gravitational phase is identical on all four branches:
$\phi_{ij}^{\mathrm{grav}} = -U_0 t/\hbar$ for all $i,j$. Therefore
\begin{equation}
|\Psi_{\mathrm{RA}}(t)\rangle
= e^{-iU_0 t/\hbar}
\frac{1}{2}\bigl(|LL\rangle + |LR\rangle + |RL\rangle
+ |RR\rangle\bigr)
= e^{-iU_0 t/\hbar}\,|{+}\rangle_A \otimes |{+}\rangle_B.
\end{equation}
The state remains separable, so $\Delta\phi_{\mathrm{grav}} = 0$.
\end{proof}

%% ===============================================================
\section{Coherence--Geometry Incompatibility}
\label{sec:incompatibility}
%% ===============================================================

The null result is not contingent on experimental parameters such as
mass, separation, or interaction time. It reflects a structural
incompatibility between the two conditions the standard BMV mechanism
requires.

\begin{theorem}[Coherence--geometry incompatibility]
\label{thm:incomp}
In RA, the following two conditions cannot simultaneously hold for
the same positional degree of freedom over the same interval:
\begin{enumerate}[nosep]
\item \emph{Branch coherence:} the alternatives $L$ and $R$ remain
in coherent superposition.
\item \emph{Branch-resolved geometry:} the alternatives $L$ and $R$
source distinct effective metrics.
\end{enumerate}
\end{theorem}

\begin{proof}
If branch coherence holds, the alternatives remain at Step~2 and
have not entered the ledger; by Lemma~\ref{lem:nonsource} they
cannot source distinct Step-5 metrics. Conversely, if branch-resolved
geometry exists, the branch label must already be represented in the
actualized ledger, meaning Step~4 has occurred and coherence is lost.
\end{proof}

This is the deeper structural reason for the null prediction. The
standard BMV mechanism requires simultaneous access to both the
coherent superposition \emph{and} the branch-dependent geometry that
would produce the entangling phase. In RA, these are architecturally
mutually exclusive.

%% ===============================================================
\section{What ``Spacetime Superposition'' Means in RA}
\label{sec:spacetime}
%% ===============================================================

The standard BMV framework implicitly represents the evolved state
as a superposition of matter--geometry composites:
\begin{equation}
\frac{1}{2}\bigl(
|LL\rangle|g_{LL}\rangle
+ |LR\rangle|g_{LR}\rangle
+ |RL\rangle|g_{RL}\rangle
+ |RR\rangle|g_{RR}\rangle
\bigr).
\label{eq:standard_superposition}
\end{equation}
RA replaces this with
\begin{equation}
\biggl(\frac{1}{2}\sum_{i,j \in \{L,R\}} |ij\rangle\biggr)
\otimes |g[\DAG_n]\rangle.
\label{eq:ra_superposition}
\end{equation}

The phrase ``spacetime superposition'' must therefore be
disambiguated:

\begin{itemize}[nosep]
\item \textbf{Forbidden in RA:} A superposition of presently actual
effective metrics---Eq.~\eqref{eq:standard_superposition}.
\item \textbf{Permitted in RA:} A superposition of Step-2
growth channels whose later actualization would yield different
Step-5 metrics---i.e., superposition over geometry-producing
\emph{futures}, not over the presently actual geometry.
\end{itemize}

This distinction is not interpretive decoration. It is the
architectural content of the Step-2/Step-5 separation and the source
of the null prediction.

%% ===============================================================
\section{Relation to Competing Frameworks}
\label{sec:comparison}
%% ===============================================================

It is important to distinguish the RA null prediction from two
superficially similar positions.

\subsection{Quantized gravity}

In frameworks where gravity is a quantized field, the metric is a
quantum degree of freedom that can be placed in superposition. Each
branch of the matter superposition sources its own branch of the
metric superposition. The standard BMV mechanism operates, and the
prediction is positive: $\Delta\phi \neq 0$.

RA disagrees. The metric is not a quantum degree of freedom. It is
the macroscopic record of actualized history.

\subsection{Semiclassical gravity}

In semiclassical gravity, the metric is sourced by the expectation
value of the stress-energy tensor:
$G_{\mu\nu} = 8\pi G\, \langle\hat{T}_{\mu\nu}\rangle$.
A spatial superposition would produce an averaged field that is
neither the $|L\rangle$ nor the $|R\rangle$ configuration but a
weighted mean. The BMV prediction in this framework is ambiguous:
some treatments predict partial entanglement, others predict none.

RA disagrees with semiclassical gravity as well. The source is not
the expectation value of unactualized branch structure. The source
is the actualized ledger alone:
\begin{equation}
G_{\mu\nu} = 8\pi G\, \Pact[T_{\mu\nu}],
\end{equation}
where $\Pact$ is the actualization projector that retains only
on-shell, ledger-written contributions.

Thus RA's null result is not ``gravity is classical.'' It is
specifically \emph{gravity is actualized}: sourced only by irreversible
causal facts, not by reversible quantum possibilities.

%% ===============================================================
\section{Experimental Consequences}
\label{sec:experiment}
%% ===============================================================

\subsection{The prediction}

In an ideal BMV experiment with all nongravitational channels
excluded, RA predicts:
\begin{equation}
\Delta\phi_{\mathrm{grav}}^{\mathrm{RA}} = 0,
\end{equation}
and any entanglement monotone attributed solely to gravity must
vanish:
\begin{equation}
\mathcal{E}_{\mathrm{grav}}^{\mathrm{RA}} = 0.
\end{equation}

\subsection{What RA allows}

RA does not predict the complete absence of all experimental
signatures. The following effects may be present:
\begin{itemize}[nosep]
\item Common-mode gravitational phases (global, nonfactorizable
contribution zero).
\item Environmental decoherence from nongravitational channels.
\item Electromagnetic or Casimir contamination.
\item Trap-dependent systematic phases.
\end{itemize}
The prediction is that after removal of these nongravitational
contributions, no purely gravitational entangling signal remains.

\subsection{Falsification condition}

A positive BMV result---specifically, a nonzero entanglement witness
attributed to gravity after exclusion of all nongravitational
channels---would directly falsify the Step-2/Step-5 separation of
the Engine of Becoming. It would demonstrate that the effective
metric responds to unactualized quantum branch structure, collapsing
the distinction between possibility and actuality that is central
to RA.

The logic is unusually clean:
\begin{itemize}[nosep]
\item \textbf{Null BMV:} consistent with RA.
\item \textbf{Positive BMV:} incompatible with present-form RA.
\end{itemize}

This makes the BMV experiment one of the sharpest near-term
empirical discriminators for the RA programme.

%% ===============================================================
\section{Discussion}
\label{sec:discussion}
%% ===============================================================

The RA-BMV null prediction reveals something deep about the framework.
RA does not merely classify gravity as classical or quantum. It makes
a more specific structural claim: \emph{gravity is the macroscopic
description of what has already become actual.}

The metric $g_{\mu\nu} = \mathcal{G}[G_n]$ is not a fundamental
field. It is the large-scale order of the actualized causal ledger.
Asking whether gravity can be ``in superposition'' is, in RA, a
category error: you cannot superpose the record of what has happened,
because what has happened is by definition actual and definite.

This does not mean RA denies quantum phenomena in gravitational
settings. Gravitational effects such as gravitational redshift,
time dilation, and frame-dragging are all present in RA as
consequences of the actualized graph's large-scale structure. What
RA denies is that the metric itself can be placed in a coherent
superposition of branch-resolved configurations, because that would
require the actualized ledger to contain mutually exclusive entries
simultaneously.

The BMV experiment is therefore not merely a test of whether gravity
is ``quantum enough'' to mediate entanglement. It is a test of
whether the effective spacetime geometry belongs to the space of
actualized facts or to the space of unresolved possibilities. RA
predicts the former. A positive BMV result would require the latter.

%% ===============================================================
\section{Conclusion}
\label{sec:conclusion}
%% ===============================================================

We have shown that Relational Actualism predicts a strict null result
for gravity-mediated entanglement in the BMV protocol. The prediction
follows from a single principle: the effective metric is sourced only
by the actualized causal ledger, not by unactualized quantum branch
structure. The coherence--geometry incompatibility theorem
(Theorem~\ref{thm:incomp}) establishes that the two conditions
required by the standard BMV mechanism---maintained positional
coherence and branch-resolved gravitational geometry---are
structurally mutually exclusive in RA.

Multiple experimental groups are currently developing BMV-type
experiments \cite{Bose2017,Marletto2017}. A null result would be
consistent with RA and would constrain frameworks in which gravity
is a quantized degree of freedom. A positive result would directly
falsify the Step-2/Step-5 separation that is the architectural
foundation of Relational Actualism.

Either outcome advances physics.

\begin{thebibliography}{99}

\bibitem{Bose2017}
S.~Bose, A.~Mazumdar, G.~W.~Morley, H.~Ulbricht, M.~Toro\v{s},
M.~Paternostro, A.~A.~Geraci, P.~F.~Barker, M.~S.~Kim, and
G.~Milburn, ``Spin entanglement witness for quantum gravity,''
\emph{Phys.\ Rev.\ Lett.}\ \textbf{119}, 240401 (2017).

\bibitem{Marletto2017}
C.~Marletto and V.~Vedral, ``Gravitationally induced
entanglement between two massive particles is sufficient
evidence of quantum effects of gravity,''
\emph{Phys.\ Rev.\ Lett.}\ \textbf{119}, 240402 (2017).

\bibitem{RAQM}
J.~F.~Sandeman, ``Relational Actualism and Quantum Mechanics,''
Zenodo (2026). \texttt{doi:10.5281/zenodo.19174380}.

\bibitem{RAGC}
J.~F.~Sandeman, ``Relational Actualism: Gravity and Cosmology,''
Zenodo (2026). \texttt{doi:10.5281/zenodo.19197999}.

\bibitem{EB}
J.~F.~Sandeman, ``The Engine of Becoming: Quantum Gravity as
Covariant Stochastic Graph Rewriting,''
Zenodo (2026). \texttt{doi:10.5281/zenodo.19198120}.

\end{thebibliography}

\end{document}
